%% main.tex — Desktop Robot / Graduated Autonomy CHI 2027 submission
%% Working title placeholder; revise once results are in.
\documentclass[sigconf,review,anonymous]{acmart}
%\documentclass[manuscript, screen, review, anonymous]{acmart}
\usepackage{subcaption}
\usepackage{placeins}

%% Keep figures/tables close to their references
\setlength{\textfloatsep}{8pt plus 2pt minus 2pt}
\setlength{\floatsep}{8pt plus 2pt minus 2pt}
\setlength{\intextsep}{8pt plus 2pt minus 2pt}
\renewcommand{\topfraction}{0.9}
\renewcommand{\bottomfraction}{0.8}
\renewcommand{\textfraction}{0.07}
\renewcommand{\floatpagefraction}{0.75}

\AtBeginDocument{%
  \providecommand\BibTeX{{%
    Bib\TeX}}}

%% Rights — fill in after CHI acceptance
\setcopyright{acmlicensed}
\copyrightyear{2027}
\acmYear{2027}
\acmDOI{XXXXXXX.XXXXXXX}
%% These commands are for a PROCEEDINGS abstract or paper.
\acmConference[HRI '27]{ACM International Conference on Human-Robot Interaction}{March 08--12,
  2027}{Santa Clara, CA}

%% Use author-year citation style for CHI
\citestyle{acmauthoryear}

\begin{document}

%% =====================================================================
%% Title — working title; finalize after main study
%% =====================================================================
\title[Tell Me Where to Go]{Tell Me Where to Go: How Users Establish Language–Action Mappings in Verbal Robot Control}

%% =====================================================================
%% Authors — anonymous in review mode; fill in for camera-ready
%% =====================================================================
\author{Anonymous Author(s)}
\affiliation{%
  \institution{Anonymous Institution}
  \country{}
}

%% --- For non-anonymous version, replace block above with the team:
%% \author{Ruixiang Han}
%% \email{rh652@cornell.edu}
%% \affiliation{\institution{Cornell Tech} \city{New York} \country{USA}}
%% \author{Sophie Bai}
%% \author{Sirui Wang}
%% \author{Xiang Chang}
%% \author{Angela Gui}
%% \author{Angelique Taylor}
%% \author{Wendy Ju}
%% \affiliation{\institution{Cornell Tech} \city{New York} \country{USA}}

%% =====================================================================
%% Abstract — write last, after Study 1 + Study 2 results
%% =====================================================================

\begin{abstract}
Recent advances in large language models and robot foundation models increasingly enable robots to interpret open-ended natural-language instructions. Yet, stronger language understanding does not guarantee that human intent can be reliably translated into embodied action. Natural language remains context-dependent, while spatial reference, control dynamics, sensing, and physical execution can introduce additional uncertainty between what a person says and what a robot ultimately does. Understanding how people navigate these language-to-action mismatches is therefore important for designing future language-enabled robots.
We study low-level verbal motion control as a controlled lens on this problem, reducing high-level planning complexity while preserving the challenge of grounding natural language in situated physical behavior. We conducted a Wizard-of-Oz study with 12 participants using a desktop robot across two complementary tasks: an open-ended, situated navigation task and a constrained path-following task involving three predefined trajectories. Through analysis of participants' verbal commands, interaction behaviors, think-aloud responses, and post-task interviews, we examine how people control, adjust, and repair verbal robot interaction.
We derive a taxonomy of verbal robot-control strategies spanning movement magnitude, temporal flow control, spatial frames of reference, and user inference about robot behavior and errors. Our findings suggest that users do not merely issue commands; they actively work to make the mapping between language and robot behavior locally interpretable and controllable. Based on these findings, we derive design implications for future voice-enabled robots that support clarification, calibration, interruption, memory, and actionable feedback. More broadly, we argue that increasingly open language interfaces do not eliminate the grounding problem, but shift part of it from predefined system design to ongoing human--robot interaction.
\end{abstract}



%% =====================================================================
%% CCS concepts + keywords — finalize close to submission
%% =====================================================================
\begin{CCSXML}
<ccs2012>
   <concept>
       <concept_id>10003120.10003121</concept_id>
       <concept_desc>Human-centered computing~Human computer interaction (HCI)</concept_desc>
       <concept_significance>500</concept_significance>
   </concept>
   <concept>
       <concept_id>10010147.10010178.10010224</concept_id>
       <concept_desc>Computing methodologies~Robotic planning</concept_desc>
       <concept_significance>300</concept_significance>
   </concept>
</ccs2012>
\end{CCSXML}

\ccsdesc[500]{Human-centered computing~Human computer interaction (HCI)}
\ccsdesc[300]{Computing methodologies~Robotic planning}

\keywords{Human--robot interaction, voice control, Wizard-of-Oz, graduated autonomy, intent ambiguity, shared agency}

\maketitle

%% =====================================================================
%% Body
%% =====================================================================
%% 01-introduction.tex
%% Status: WRITTEN (v1) — framing stable; expect light revision after results.
%% Source: NL_Robot_Synthesis_Short Part 1 + Part 3; Summary - Graduated Autonomy


\section{Introduction}

Recent advances in large language models (LLMs) and robot foundation models have substantially expanded robots' ability to interpret open-ended natural-language instructions, moving beyond predefined command vocabularies and rigid interaction grammars. \citep{ichter2023saycan, driess2023palme, brohan2023rt2} This shift opens a broader interaction space in which people can express goals, preferences, corrections, and situated needs without first learning a machine-specified control language.

Yet, the ability to interpret increasingly flexible language does not mean that language can be unambiguously translated into embodied behavior. Natural language is inherently context-dependent: an instruction such as ``move a little closer'' leaves the magnitude of ``a little'' underspecified, while expressions such as ``left'' may depend on whether the speaker adopts the human's or the robot's frame of reference. \citep{kocher2020left, li2016spatial} Importantly, uncertainty does not end once the intended meaning has been correctly interpreted. Even an apparently precise instruction such as ``turn 30 degrees'' may result in different physical outcomes due to sensing, control, friction, or other properties of the embodied system. Thus, the path from what a person intends to what a robot ultimately does is not a simple one-to-one mapping.

We conceptualize this process as a grounding chain from \textit{human intent} to \textit{linguistic expression}, \textit{robot interpretation}, and ultimately \textit{embodied execution}. Prior work on natural-language robot interaction has extensively investigated how robots ground linguistic expressions to objects, spatial relations, trajectories, and actions. \citep{tellex2011commands, kollar2010directions, matuszek2013parse} As language-enabled robots become increasingly capable, however, an equally important interactional question emerges: what happens on the human side when the mapping between language and embodied action is incomplete, uncertain, or inconsistent? How do users determine what their words mean for a particular robot, recognize when that mapping breaks down, and adjust the interaction accordingly?

These processes can be difficult to isolate in complex, high-level tasks. If a robot fails when asked to ``organize the desk,'' for example, the source of failure may lie in goal interpretation, task planning, perception, navigation, manipulation, or physical execution. We therefore study verbal motion control as a more controlled lens on language-to-action grounding. Commands such as ``move forward,'' ``turn left,'' or ``keep going until I say stop'' reduce high-level planning complexity while preserving the fundamental challenge of translating natural language into situated physical action. \citep{marge2010routes, marge2015recovery, fong2003questions}

Accordingly, we ask:

\begin{itemize}
    \item \textbf{RQ1:} What verbal and interaction strategies do people use when controlling a robot's physical motion through natural language?
    \item \textbf{RQ2:} What challenges and interaction needs emerge as people use, adjust, and repair verbal robot control across different task demands?
\end{itemize}

To answer these questions, we conducted a Wizard-of-Oz study with 12 participants using a desktop robot. Participants completed two complementary verbal-control tasks: an open-ended, situated navigation task that allowed them to determine how to reach a meaningful goal, and a constrained path-following task that required more precise adherence to three predefined trajectories. Together, these tasks allowed us to observe verbal control under both flexible and precision-sensitive conditions. We analyzed participants' verbal commands, interaction behaviors, think-aloud responses, and post-task interviews.

Our work makes three contributions:

\begin{itemize}
    \item We present an empirically derived \textbf{taxonomy of verbal robot-control strategies}, characterizing how participants expressed movement magnitude, managed the temporal flow of action, handled spatial frames of reference, and reasoned about robot behavior and errors.
    
    \item We characterize the \textbf{interaction challenges and user responses} that arise when natural-language instructions must be translated into situated physical action, including command--execution mismatch, interaction effort, feedback needs, and strategies for adjusting or repairing control.
    
    \item Building on these findings, we derive \textbf{design implications for future language-enabled robots}, arguing for systems that not only interpret user commands but also support clarification, calibration, interruption, memory, and actionable feedback as part of ongoing human--robot interaction.
\end{itemize}

%% 02-related_work.tex
%% Status: WRITTEN (v1) — does not depend on results.
%% Citations cross-checked against references.bib.

\section{Related Work}
\label{sec:related-work}


Verbal robot control requires people and robots to coordinate both the meaning of an instruction and the action it produces. We draw on research in language grounding, conversational conventions, spatial reference, shared control, and communication repair to examine how these forms of coordination develop through interaction.

\subsection{Grounding Language in Action}

Research on natural language interfaces has developed several approaches to connecting instructions with robot behavior. Generalized Grounding Graphs associate the compositional structure of language with objects, places, paths, and events to infer action plans \citep{tellex2011commands}. \citet{matuszek2013parse} developed a parser that translates route instructions into a procedural language representing actions and control structures. More recently, SayCan combined language model scores with value functions for available skills, allowing the system to consider both the relevance of an action and its feasibility for a particular robot \citep{ichter2023saycan}. Across these approaches, interpretation depends on the relationship between language, representations of the environment, and the actions available to the system.

Studies of command elicitation examine how people formulate the instructions that these systems must interpret. \citet{marge2010routes} collected spoken route descriptions across schematic, virtual, and physical scenes. Participants preferred metric descriptions but also used landmarks to specify movement. Because the robots remained stationary and participants described complete routes, the study primarily concerns instruction formulation. \citet{kocher2020left} examined a more interactive setting in which 38 children aged four to nine directed a virtual robot controlled by an experimenter, as well as their caregivers. Children used precise and vague distances, corrective feedback, and strategies that involved turning until told to stop. The experimenter assigned a movement amount to an initially unspecified distance expression and scaled later uses accordingly. These observations provide a close precedent for local units and ongoing correction, although the participants and virtual setting differ from those of the present study.

The SCOUT corpus extends this work to sustained dialogue during collaborative exploration. It combines multiple Wizard-of-Oz studies in which participants gave spoken instructions to a remotely located robot, with annotations connecting utterances to intentions and dialogue structure \citep{lukin2024scout}. Such corpora make it possible to examine instructions in relation to preceding exchanges and information about the environment. Our study builds on this tradition by considering movement magnitude, action timing, spatial reference, and interpretations of failure within a common account of verbal control across two physical navigation tasks.

\subsection{Local Conventions and Spatial Reference}

Grounding language in a computational representation is one part of establishing understanding between interaction partners. \citet{clark1991grounding} describe conversational grounding as a collaborative process through which participants establish understanding sufficient for their current purposes. Acknowledgments, subsequent contributions, and other evidence allow partners to assess this understanding, while the communication medium shapes the effort involved. \citet{brennan1996pacts} show how people develop provisional conceptualizations for referring to objects with particular partners. These conceptual pacts explain why the interpretation of an expression depends on a shared interaction history as well as its descriptive content.

Research in HRI documents several forms of linguistic adaptation. In a collaborative map task, \citet{brandstetter2017entrainment} found that a robot's choice of words influenced participants' later naming, including after the interaction had ended. \citet{asano2022alignment} found greater lexical alignment when students interacted individually with a teachable robot than in the portions of triadic interactions directed toward the robot. The relationship between alignment and rapport was more complex than predicted, suggesting that vocabulary convergence should be interpreted in relation to the interaction in which it occurs. A recent preprint by \citet{zhang2026alignment} organizes conceptual alignment according to its triggers and the levels of understanding involved. Its taxonomy was developed from dialogue between people performing concept-sorting tasks, so its application to physical robot control remains an analytical extension. Together, these accounts motivate studying local movement expressions as possible conventions whose shared meaning must be established through use.

Spatial reference introduces a related challenge because partners may describe the same environment from different perspectives. \citet{li2016spatial} found that written instructions for tabletop object selection frequently relied on a perspective that was not explicitly stated. Perspective and the complexity of spatial relations were associated with difficulties in interpreting those instructions. Although these studies examined written descriptions and subsequent object selection, they demonstrate the importance of establishing which frame a spatial expression invokes. Physical motion adds an opportunity to test that interpretation through its consequences. We examine how users respond to uncertainty about reference frames through clarification and exploratory movement, alongside their choices about the extent and timing of an action.

\subsection{Online Correction and Control Authority}

The meaning of a correction depends partly on what the robot is doing when it is issued. \citet{raux2010corrections} distinguished corrections of omitted actions, incorrect actions, and action degree, examining their linguistic, prosodic, and temporal characteristics in dialogue between people interacting with a simulated robot clerk. Their analysis connects corrective speech to the progress of an action. Dialogue can also determine who takes responsibility for resolving a problem. In collaborative teleoperation, \citet{fong2003questions} developed a system in which robots requested human assistance with perception and cognition while treating human attention as a limited resource. Although this system used a PDA interface, it established dialogue as part of the control loop. \citet{nikolaidis2017mutual} subsequently modeled human adaptability in shared autonomy, enabling a robot to adjust its behavior according to a user's willingness to change strategies.

Recent systems integrate language into ongoing control in different ways. LILAC uses language corrections during manipulation to update a learned control space that the user operates through a joystick \citep{cui2023lilac}. Language therefore changes how manual input guides the robot. BRIDGE allows users to modify planned position, velocity, and force through speech during physical assistance, while the robot provides assurances or asks for clarification \citep{wang2026bridge}. A study with 18 older adults demonstrated these adjustments across assistive tasks. The implementation pauses motion after a completed utterance while computing a modification, making the timing of an intervention an explicit feature of the interaction. \citet{cao2025interruption} addressed another aspect of timing by implementing conversational interruption handling that responds to the interrupter's intent. Their evaluation concerned discussion tasks and therefore provides evidence about managing conversational turns, rather than physical stopping performance.

This literature establishes correction and the allocation of control as central concerns in HRI. We examine how users express control arrangements in speech, including bounded movements, motion that continues until the user intervenes, repeated authorization to proceed, and delegated sequences of actions. These arrangements specify responsibility for continuation and termination alongside the intended movement. Relating them to magnitude and spatial reference helps identify what an interface must make explicit for verbal control to remain understandable and manageable. Comparing the efficiency and safety of these arrangements would require further evaluation.

\subsection{Repair, Transparency, and Calibration}

Repair requires an account of what went wrong. \citet{marge2015recovery} distinguish ambiguity about a referent from instructions that cannot be executed in the current environment. Their crowdsourced study found that recovery responses used relevant visual information to resolve ambiguity and proposed alternative plans when a request was infeasible. \citet{chai2014ground} examined how robot feedback supports common ground in situated dialogue. They found that minimal acceptance feedback could leave people believing that understanding had been established when mutual understanding was incomplete. These studies show why a successful acknowledgment cannot be assumed to identify either the robot's interpretation or the cause of a subsequent mismatch.

Providing access to intermediate processing can help users assess that interpretation. \citet{perlmutter2016transparency} exposed perception and inference outputs in a system for understanding situated language. In a study with 20 participants, visualizations on a screen improved selected measures of command communication, but the benefits did not reliably persist after the displays were removed. The findings distinguish support available during an interaction from lasting changes in a user's mental model. They also suggest that the value of transparency lies partly in helping people decide how to formulate or revise an instruction.

Recent work with language models offers further mechanisms for clarification and personalization. KnowNo uses calibrated prediction sets to determine when a planner should seek human help with an ambiguous choice \citep{ren2023knowno}. TidyBot generalizes object-placement preferences from user examples to support personalized tidying \citep{wu2023tidybot}. These capabilities provide precedents for adapting robot behavior to information supplied by users. Applying them to verbal motion control would additionally require mechanisms for establishing and revising movement units, spatial frames, and stopping arrangements in relation to the robot's observed behavior.

Our study connects these concerns by examining how users specify and regulate motion while reasoning about mismatches and possible repairs. The taxonomy relates choices about magnitude, timing, and spatial reference to users' interpretations of robot behavior. It motivates feedback that distinguishes what the system received, how it interpreted the instruction, and what it executed. We use the idea of a local control convention to interpret users' efforts to establish workable mappings within an interaction. Whether such conventions become stable, or whether particular calibration mechanisms improve control, remains a question for subsequent evaluation.

%% 03-framework.tex — Conceptual Framework
%% Status: WRITTEN (v1) — concepts stable; expect light revision after results.
%% Source: NL_Robot_Synthesis_Short Part 1; Summary - Graduated Autonomy

\section{Method}
\label{sec:method}
\subsection{Desktop Robot}
We created a differential-drive desktop robot that runs ROS 2 Jazzy on a Raspberry Pi 5 with Dynamixel XL330 motors. It adapts a microphone (ReSpeaker Mic Array v2.0), and a plug-in speaker as the output. These are the only user-level communication channel. We intentional keep the feature of the robot minimal which can only adapt vocie input/output to comply with our study design purpose. Meanwhile, the embodiment is intentionally designed as round, which means we provide ambiguous toward the orientation, so participants should be able to communicate and negotiate with the robot.
(instruction: put desktop robot picture here)
\subsection{Deployment and Setting}
We recruited 12 participants from our campus facility. We conducted a disclosed wizard of Oz experiment with them, where they were told the natrue of the study when they sign the consent form. We request them to communicate with the robot directly, providing only the low level instruction. When they try to provide high level instruction such as "go toward me", or "go toward the vessel", the robot will provide. The finish task one followed by task two.
\subsubsection{Task one: exploration}
Task one composites a everday life desktop setting: laptop, vassel, cup, and other office tools were scattered on the desktop. The desktop robot sits at the other corner opposite of the participants, and the participants will need to use verbal instruction to make the robot to move toward them and meanwhile circumstance all the obstacles along the way.  
(instruction: Put study1 setting, and study one participants picture here)
\subsubsection{Task two: projected trajectory}
Task two requests participant to order the robot to move along the pre-defined trajectory and move toward themselves. Instead of the open-ended first task where participants have some improvisation space, this task require a stricter, concise level of wording description for participant to navigate the robot along the line. To give participant time to get familiar themselves with the task, the task is therefore designed in ascending order of difficulty. The first task is a simple right-angle turn, the second is the zig-zag line, and the final and thrid one is a complex route. 
(instruction: Put the study2 setting1-3 as pictures in here. Together)
\subsection{Video analysis}
Three researchers independly watched the video and come out with the coding scheme of the strategies that participants used for the interaction with the desktop robot. Then we gather together to finalized the coding that we are going to share. We created a browser-based video annotator, and each researcher independly coded these video, and then we output these file as 
(instruction: Put web interface here)


\section{Result}

\subsection{Verbal and Interaction Strategies for Controlling Robot Motion}
For \textbf{RQ1}, we seek to characterize the verbal and interaction strategies participants use when controlling a robot's physical motion through natural language. Across the two tasks, participants varied not only in how they described motion, but also in how they controlled its temporal progression, established spatial reference frames, and reasoned about the robot's behavior. We organized these behaviors into three dimensions: movement command specification, command sequencing, and frame-of-reference management.

\subsubsection{Movement Command Specification}
Participants drew on five strategies to specify robot motion, often shifting between or combining them during interaction: quantitative measures, qualitative expressions, user-defined units, relative adjustments, and environmental references. \textbf{Quantitative measures} involved specifying movement using explicit magnitude such as centimeters, degrees, or time (P3, P4, P5, P8, P10). Participants gave commands such as \emph{``move forward 50 centimeters''} and \emph{``turn left 30 degrees.''} Time was also used as a proxy for distance, as in \emph{``walk straight for 5 seconds.''} In contrast, \textbf{qualitative expressions} relied on vague magnitude terms such as \emph{``a little bit,''} \emph{``a larger step,''} or \emph{``a really long time''} (P6, P7, P8, P9, P12). These expressions were often refined incrementally after participants observed the robot's movement, for example by asking it to move \emph{``a little bit more.''}

Some participants established their own movement scales through \textbf{user-defined units} (P5, P7, P9, P11). P5 explicitly defined a particular movement amount as \emph{``one unit,''} while others used measures such as \emph{``steps''} or wheel rotations without formally specifying their physical magnitude. Participants also used \textbf{relative adjustments} to specify new movements in relation to previous ones (P4, P5, P8, P9, P10, P12). Rather than providing a new absolute magnitude, they requested the same amount, \emph{``twice the amount,''} \emph{``5 more degrees,''} or simply instructed the robot to \emph{``keep going.''} Finally, we observed one participants used \textbf{environmental references} to describe desired movement in relation to objects or destinations rather than low-level motion parameters. For example, P1 instructed the robot to \emph{``go in between a frame and a cup.''} Because the robot could not visually ground these environmental references, participants quickly shifted away from this strategy toward primitive movement commands.

\subsubsection{Movement Command Sequencing}
Participants used three strategies to organize commands and regulate how robot motion unfolded over time: discrete step-and-stop control, user-terminated continuous commands (UTCs), and command batching. In \textbf{discrete step-and-stop control}, participants (P1, P2, P3, P6) issued short movement primitives, observed the result, and stopped the robot before continuing. Typical sequences included \emph{``Turn right... Stop''} and \emph{``Forward... Stop.''} Stop commands often marked either a movement boundary or the detection of an error, and were sometimes accompanied by repair expressions such as \emph{``No! Turn left''} or \emph{``Sorry, stop.''}

In \textbf{user-terminated continuous commands (UTCs)}, participants allowed the robot to continue moving until they later issued a stop. Unlike step-and-stop control, the terminal stop was generally anticipated as part of the control strategy rather than introduced primarily as a correction. Some participants (P1, P2, P4, P8, P9, P11) explicitly established this stopping arrangement in advance, as in \emph{``Keep moving forward until I say stop.''} Others (P9, P10, P11, P12) sustained the ongoing movement through repeated continuation cues, such as \emph{``Keep going. Keep going... Stop.''} Finally, in \textbf{command batching}, participants (P3, P5, P7, P8, P12) combined multiple movement primitives within a single utterance. Rather than issuing each instruction separately, they specified a sequence of actions, such as moving forward, rotating, and then moving forward again.

\subsubsection{Managing Frame-of-Reference Ambiguity}
Participants used several strategies to resolve ambiguity in directional commands such as left, right, forward, and backward, which could be interpreted from the participant's perspective, the robot's perspective, or another external frame of reference. Some participants \textbf{anchored directional commands} to the robot's body frame. P10 and P12 explicitly referenced the robot's perspective using expressions such as \emph{``Turn to your right a little bit''} or \emph{``Turn a little bit to your left.''} P11 instead used \textbf{orientation-independent terms}, such as \emph{``clockwise''} and \emph{``counterclockwise,''} to avoid left--right ambiguity. P11 also distinguished rotation from translation, using \emph{``spin''} to indicate rotation in place and \emph{``move''} or \emph{``go''} to indicate translation, as in \emph{``Stay in the same position; spin 45 degrees.''}

Other participants resolved uncertainty interactively. P6, P7, and P8 \textbf{asked about the robot's orientation} before continuing, using questions such as \emph{``Which way is your head?''}, \emph{``Which side is the front?''}, and \emph{``Wait, which way is forward?''} P7 described the ambiguity directly: \emph{``I don't know whether I should tell the direction according to my view or the robot's view.''} Participants also used \textbf{test-and-recalibrate} strategies when the reference frame remained unclear (P2, P8, P11). P8 tested forward and backward movements, while P11 confirmed the observed mapping with \emph{``Okay, so that's forwards.''} P2 instead reset the robot to a canonical orientation using \emph{``Turn to the front''} before continuing.

\subsection{Challenges and Interaction Needs}
For RQ2, we seek to learn about the challenges participants encountered as they adjusted and repaired verbal robot control, as well as the interaction support they expected across different task demands. Participants generally found the free navigation task more tolerant of imprecision because they could adapt the route as long as the robot reached the intended destination (P3, P8, P9, P10, P11), whereas the constrained path task following required finer-grained turns, closer trajectory adherence, and more frequent corrections (P5, P8, P11). Participants also reported mixed overall experiences: some found the interaction comfortable or manageable when the robot generally followed their instructions (P1, P4, P5, P7), while others described it as frustrating or tiring because of repeated corrections and the number of low-level commands required (P2, P3, P8, P11). These differences provide context for the challenges and interaction needs described below.

\subsubsection{Control and Execution Challenges}
The most common challenge was a mismatch between the motion participants specified and the motion the robot produced (P3, P4, P5, P7, P8, P9). This mismatch occurred even when participants used explicit numerical commands. Rotation was particularly difficult to predict (P3, P4, P9). P3 explained that \emph{``sometimes when I say maybe 30 degrees, it takes like 60 degrees or just like 90 degrees,''} while P9 observed that \emph{``every time the angle that it makes turns is kind of different.''} P4 similarly noted that \emph{``what I thought was not actually what the robot was responding to.''} Participants encountered similar problems with distance, speed, and straight-line movement (P5, P7, P8, P9). P5 reported that \emph{``when the robot moved forward, it was not exactly moving the same amount as I wanted it to.''} P7 had to \emph{``guess how many steps it takes''} and then adjust subsequent commands, while P8 repeatedly corrected the robot because it drifted slightly to the left or right rather than maintaining a straight trajectory.

Participants also struggled to translate spatial intentions into verbal quantities. P10 stated, \emph{``It is hard for me to estimate how long it should walk, as well as what exact degree it should rotate.''} As a result, participants sometimes relied on approximations, such as specifying movement duration rather than physical distance. These difficulties increased the interaction effort required to control the robot. P2 described the process as \emph{``tedious''} because every small movement required another instruction, while P8 recalled having to issue a forward command approximately 20 times, which was \emph{``tiring.''} Repeated micro-instructions therefore made verbal control burdensome, particularly when precise positioning required frequent observation and correction.

\subsubsection{Challenges in Diagnosing and Repairing Failures}
A further challenge was determining why the robot had behaved unexpectedly and how to recover from the resulting mismatch. Participants did not always know whether a failure originated in their command, the robot's interpretation, or its physical execution. Some attributed problems to the robot's interpretation of their language. P6, for example, stated, \emph{``I feel like maybe the numerics are not working.''} Others attributed discrepancies to physical execution. After observing an inaccurate turn, P11 reasoned, \emph{``It turned 170 because of friction, so let's try 190,''} and compensated by modifying the next command. Participants also considered whether their own wording had contributed to the problem. P10 reflected, \emph{``Maybe I'm not using the best command to communicate.''}

Because the source of failure was often uncertain, repair required participants to experiment with different responses. They reformulated commands, adjusted numerical parameters, or replaced an ongoing instruction with a new one. P4 described issuing a follow-up instruction to \emph{``break the old one and continue this new routine,''} while P5 preferred asking the robot to \emph{``repeat the last step.''} In some cases, participants also disregarded robot feedback when it conflicted with what they could directly observe. P1 instructed the robot to \emph{``Go forward, trust me,''} while P11 responded to an obstacle warning with \emph{``That's a cup, you don't have to go around.''} The effort required to diagnose and recover from these failures also shaped participants' expectations for how the robot should support the interaction.

\subsubsection{Interaction Needs}
Participants wanted the robot to provide enough feedback to make its interpretation and execution understandable. They expected it to acknowledge commands, state the action it intended to perform, and report completion or progress. P4 suggested brief acknowledgments such as \emph{``okay''} or \emph{``got it,''} while P9 wanted the robot to state, for example, \emph{``I am now going to turn right for 60 degrees,''} and confirm when the action was complete. Such feedback was particularly important for distinguishing whether an unexpected outcome resulted from misunderstanding or physical execution. Without feedback, P10 assumed that the robot \emph{``does not understand me,''} while P11 wanted to know whether a problem came from \emph{``a faulty robot''} or the \emph{``command.''} Participants also valued obstacle feedback (P2, P3, P7, P9), although its usefulness depended on what they could already observe. At the same time, they did not want excessive verbal communication. P1 stated that the robot should \emph{``talk back when necessary but shouldn't be talking back too much,''} and P6 suggested that routine information such as obstacle detection could instead be communicated through a non-verbal signal, such as a red light.

When commands were ambiguous, participants expected the robot to ask for clarification rather than silently select an interpretation. However, they also expected the robot to retain the resulting clarification. P8 noted that repeatedly asking what \emph{``a little bit''} meant would become \emph{``kind of annoying.''} Participants therefore favored clarification combined with memory, so that locally established meanings could be reused during the interaction. Participants further expected the robot to disclose its capabilities and limitations and to support calibration before or during the task. P4 suggested providing example command forms and a brief calibration procedure to reduce trial-and-error interaction. More generally, participants expected negotiation to remain proportional to the task. P7 stated that the robot should generally \emph{``strictly follow the user's instructions,''} while explaining its reasoning when it could not. P11 similarly emphasized that additional interaction effort should feel justified, stating, \emph{``I want my efforts to be necessary.''}

\section{Discussion}

\subsection{From Issuing Commands to Stabilizing Language--Action Mappings}

Our findings suggest that verbal robot control involves more than identifying a ``correct'' way to phrase an instruction. Across different strategies, participants sought to make the relationship between what they said and what the robot did sufficiently stable for the interaction to continue. Prior work has shown that conversational partners develop locally shared expressions and that spatial language depends on negotiated or inferred frames of reference. \citep{brennan1996pacts, li2016spatial} Our findings extend this perspective into embodied control. Participants moved between quantitative, qualitative, and relative expressions, introduced locally meaningful units, tested spatial mappings, changed how they controlled continuation and stopping, and reasoned about whether mismatches originated in their own expressions or in robot execution.

Importantly, this uncertainty did not reside in natural language alone. A numerically precise command could still produce an imprecise physical outcome, making linguistic precision distinct from embodied precision. Viewed together, the four dimensions of our taxonomy characterize different aspects of the mapping that users sought to stabilize: magnitude establishes how much, flow control establishes when action continues or stops, spatial reference establishes from whose frame, and error inference helps determine what should be repaired when the mapping fails. Rather than treating these behaviors as isolated command styles, we interpret them as components of a \textit{session-local control convention} through which users make verbal control locally interpretable and manageable.

\subsection{Designing Robots to Participate in Calibration}

Existing work on adaptive and shared human--robot interaction has emphasized that effective collaboration can involve adjustment across both human and robot behavior. \citep{nikolaidis2017mutual, cui2023lilac} In our study, however, much of this adjustment remained on the human side. Participants changed their wording, compensated for observed motion, selected different reference frames, adjusted the granularity of control, and inferred why an interaction had failed.

This shifts the design opportunity from simply building robots that can understand more language toward building robots that can \textit{participate in establishing what language means in the current interaction}. When an expression such as ``a little'' is underspecified, a robot can request clarification rather than silently committing to one interpretation. When users establish a useful local convention, such as a personalized movement unit, the system can retain and revise it. When spatial frames are uncertain, the robot can make its interpretation explicit; when action deviates from the intended outcome, feedback can help users distinguish interpretation errors from execution errors and select an appropriate repair. Recent LLM-enabled robotic systems already demonstrate capabilities such as uncertainty-triggered clarification and learning or retaining user-specific information. \citep{ren2023knowno, wu2023tidybot} These capabilities create an opportunity for verbal control to become less dependent on unilateral human adaptation and more actively negotiated between human and robot.

\subsection{More Capable Robots Make Interaction-Time Grounding More Important, Not Less}

Our findings also help explain why low-level verbal control remains relevant as robots become increasingly capable. Traditional robot interfaces often reduce uncertainty by constraining users to predefined commands, parameters, or interaction rules; much of the mapping between input and action is therefore established by designers before interaction begins. \citep{fong2003questions, marge2010routes} In contrast, language models increasingly allow robots to accept flexible, open-ended, and context-sensitive instructions. \citep{ichter2023saycan, driess2023palme, brohan2023rt2}

This increased flexibility does not eliminate the need for a control protocol. Instead, it shifts part of that protocol \textit{from design time to interaction time}. Expressions such as ``move a little closer,'' ``keep going until I say stop,'' or ``do it like before'' cannot be fully specified independently of the particular user, robot, physical environment, and interaction history. Low-level motion control provides a controlled setting in which this shift becomes especially visible, but the underlying challenge extends beyond navigation. As robots take on richer and more autonomous roles, they will increasingly need to establish, maintain, and repair situated conventions with people rather than merely execute mappings fixed in advance. In this sense, greater language capability does not remove the grounding problem; \textit{it changes where grounding happens and who participates in resolving it}.


\input{sections/06-limitation}
%% 08-conclusion.tex
%% Status: BLOCKED until Study 1 + Study 2 results are in.
%% Target length: ~0.5 page

\section{Conclusion}
\label{sec:conclusion}

% Recap the three contributions:
%   (1) Conceptual: Graduated Autonomy + Perceived Controllability
%   (2) Empirical: thematic findings + quantitative effects
%   (3) System: multi-agent intent-routing architecture
%
% End with the one-sentence takeaway from Summary - Graduated Autonomy:
%   "Autonomy should not be pre-defined — it should be negotiated, learned,
%    and gradually constructed through interaction."
% TODO


%% =====================================================================
%% Acknowledgments — keep commented during anonymous review
%% =====================================================================
%\begin{acks}
%  TODO
%\end{acks}

%% =====================================================================
%% References
%% =====================================================================
\bibliographystyle{ACM-Reference-Format}
\bibliography{references}

%% =====================================================================
%% Appendices (optional)
%% =====================================================================
%\appendix
%\input{sections/appendix_protocol}
%\input{sections/appendix_measures}

\end{document}
